% ============================================================================
%  附录 A：证据层级与核验状态
% ============================================================================
\section{证据层级与核验状态}
\label{app:log}

正文中的结论分为三类：带原始文献与编号的已知定理、标注为「个人复现」的自算内容、
以及仅保留在方法边界中的未推广模式。下表把正文用到的每个来源与其证据强度对齐；
研究过程的日志与附件访问记录不进入正文，仅保留在本节末尾。

\begin{table}[!htbp]
\centering
\footnotesize
\renewcommand{\arraystretch}{1.22}
\caption{来源与核验状态}
\label{tab:status}
\begin{tabularx}{\textwidth}{@{}>{\raggedright\arraybackslash}p{3.5cm}
    >{\raggedright\arraybackslash}p{2.9cm}
    >{\raggedright\arraybackslash}p{2.6cm}
    >{\raggedright\arraybackslash}X@{}}
\toprule
来源 & 状态 & 支撑正文 & 说明 \\
\midrule
Hsiang--Szczarba（1964） & 原文已核验 & 第 2.3 节 &
定理 1.1 的稳定切丛公式，与局部图像证明交叉核对 \\
Lam（1975） & 原文已核验 & 第 2.3 节 & 定理 1.1 的 Grassmann 情形 \\
Oproiu（1977） & 原文已核验 & 第 4.1、4.2、5.9 节 &
定理 1--2、公式 (1)--(6)、引理 4 及第 181--184 页的主要步骤 \\
Hiller--Stong（1981） & 原文已核验 & 第 4.3、5.3 节 &
命题 4.1--4.3 与未编号 \emph{Corollary}；对象是 immersion \\
Hiller（1980） & 原文已核验 & 第 5.6、5.7 节 &
Theorem 1.1、\textsection2 的旗流形计算、附录递归；定位为计算来源 \\
Stong（1982） & 原文已核验 & 第 5.8 节 &
旗流形评价、$w_1$ 的高度与 $G_3$ 杯积；不产生嵌入结论 \\
Milnor--Stasheff（1974） & 相关判据已核验 & 第 3.3 节 &
Whitney 对偶与 Thom 类论证（\textsection11） \\
$G_{2,2}$、$G_{2,3}$、$G_{3,6}$ 的计算 & 个人复现 & 第 5.2、5.7、5.9 节 &
商环关系、Gr\"obner 基与有限域秩核验，均已在真实商环中完成 \\
\bottomrule
\end{tabularx}
\end{table}

\FloatBarrier

\paragraph{研究管理记录。}
本文经历两轮整稿修订。第一轮修复法丛与 Thom 类论证、Oproiu 两族参数与幂次、
Hiller--Stong 的浸入界、Stong 的包含关系以及 $G_{3,6}$ 的商环非零性；第二轮统一
摘要、结论、参考文献、章节交叉引用与证据层级，并把过程性待办移出正文。复写为
\LaTeX{} 工程时另做了记号统一：旗流形投影一律记 $\pi$，$p=s-n$ 与 $m=2s-n$
分指两族的参数，正交补丛自身的类一律记 $v_j$。
